\documentclass[11pt,a4paper]{article}

\usepackage[spanish,es-noshorthands]{babel}
\usepackage[utf8]{inputenc}
\usepackage[T1]{fontenc}
\usepackage{lmodern}
\usepackage{amsmath,amssymb,amsthm}
\usepackage[margin=2.5cm]{geometry}
\usepackage{booktabs}
\usepackage{enumitem}
\usepackage[hidelinks]{hyperref}
\usepackage{tikz}
\usetikzlibrary{positioning,arrows.meta,fit,backgrounds,patterns,decorations.pathreplacing,calc}
\usepackage{algorithm}
\usepackage{algpseudocode}
\floatname{algorithm}{Algoritmo}
\algrenewcommand\algorithmicfor{\textbf{para}}
\algrenewcommand\algorithmicdo{\textbf{hacer}}
\algrenewcommand\algorithmicwhile{\textbf{mientras}}
\algrenewcommand\algorithmicif{\textbf{si}}
\algrenewcommand\algorithmicthen{\textbf{entonces}}
\algrenewcommand\algorithmicelse{\textbf{si no}}
\algrenewcommand\algorithmicend{\textbf{fin}}
\algrenewcommand\algorithmicrequire{\textbf{Entrada:}}
\algrenewcommand\algorithmicensure{\textbf{Salida:}}
\algrenewcommand\algorithmiccomment[1]{\hfill\(\triangleright\) #1}

\newtheorem{definicion}{Definición}
\newcommand{\Dias}{\mathcal{D}}
\newcommand{\Sem}{\mathcal{S}}
\newcommand{\Bases}{\mathcal{B}}
\newcommand{\Lin}{\mathcal{L}}
\newcommand{\Trab}{\mathcal{W}}
\newcommand{\Tipos}{\mathcal{T}}
\newcommand{\Franjas}{\Phi}
\newcommand{\Desc}{\mathcal{R}}

\title{Generador de cuadrantes anuales\\[0.3em]
       \large Modelo matemático y algoritmo}
\author{}
\date{Octubre de 2026}

\begin{document}
\maketitle

% =========================================================================== %
\section{Resumen}
% =========================================================================== %

El generador construye el \emph{cuadrante anual} de una zona de transporte sanitario no urgente:
para cada trabajador y cada día del año, qué turno hace o qué descanso tiene. Parte de tres
fuentes de datos: el \emph{plan funcional} (las líneas que hay que cubrir, con sus horarios y los
días en que operan), la \emph{plantilla} (quién trabaja, con sus vacaciones y condiciones) y el
\emph{calendario} de festivos de cada base.

Un cuadrante válido tiene que cumplir el convenio colectivo (descanso mínimo entre jornadas,
tope de días y de horas por semana, jornada anual) y las reglas de organización de la zona
(descanso semanal y fines de semana). Entre todos los
cuadrantes válidos se busca, por este orden de prioridad:
\begin{enumerate}[nosep]
  \item La máxima \textbf{cobertura} de las plazas del plan funcional
  \item Que cada trabajador quede lo más cerca posible de su \textbf{jornada anual}, para optimizar su uso
  \item Un reparto \textbf{equitativo} de los turnos duros (noches, tardes, partidos) y de los
        fines de semana
  \item Conseguir una \textbf{rotación} reconocible: semanas estables, sin repetir seguidas las mismas
        semanas duras.
\end{enumerate}

El problema completo, planteado como un único modelo de optimización, es demasiado grande para
resolverse con garantías en un tiempo razonable, ya fue intentado al comienzo, y, además, sería difícil de explicar. Por eso pensé en
\textbf{descomponer en pasos}, cada uno de los cuales parte del cuadrante que deja el anterior y lo
completa:
\begin{itemize}[nosep]
  \item Por un lado los pasos \emph{deterministas} colocan lo que ya viene decidido o lo que se reparte bien
        por turnos (los patrones pactados, sus cesiones, los fines de semana y festivo)
  \item Los pasos de \emph{optimización} (modelos de programación con restricciones, CP-SAT)
        deciden lo que de verdad necesita una visión global del año: las semanas de lunes a
        viernes de la plantilla fija y de los correturnos. Aunque se ha considerado usar una version greedy.
\end{itemize}
Cada paso respeta lo que dejaron los anteriores como constante, y todos comparten la misma
definición de las restricciones legales.

% =========================================================================== %
\section{El problema y la notación}
% =========================================================================== %

\subsection{El calendario}

Sea $\Dias$ el conjunto de días del año y $\Sem$ el de sus semanas ISO (de lunes a domingo) para
cada día $d$, $s(d) \in \Sem$ es su semana. Sea $\Bases$ el conjunto de \emph{bases} (los
municipios desde los que se opera). Cada base $b$ tiene su conjunto de festivos
$F_b \subseteq \Dias$, que incluye los nacionales y los de su calendario local.

\begin{definicion}[Tipo de día]
Para una base $b$, el tipo del día $d$ es
\[
  \tau_b(d) =
  \begin{cases}
    \mathrm{FES} & \text{si } d \in F_b \text{ es festivo}\\
    \mathrm{LV}  & \text{si } d \notin F_b \text{ y } d \text{ es de lunes a viernes},\\
    \mathrm{SAB} & \text{si } d \notin F_b \text{ y } d \text{ es sábado},\\
    \mathrm{DOM} & \text{si } d \notin F_b \text{ y } d \text{ es domingo},
  \end{cases}
\]
con $\Tipos = \{\mathrm{LV}, \mathrm{SAB}, \mathrm{DOM}, \mathrm{FES}\}$. El festivo manda sobre
el día de la semana: un sábado festivo es de tipo $\mathrm{FES}$.
\end{definicion}

\subsection{Las líneas}

Una \emph{línea} $l \in \Lin$ es un turno del plan funcional que se repite todos los días en que
opera. Cada línea tiene:
\begin{itemize}[nosep]
  \item Su base $b(l) \in \Bases$ y los tipos de día en que opera, $O(l) \subseteq \Tipos$
  \item Su horario: hora de entrada y de salida. Para un día $d$, su intervalo real es
        $[\,\mathrm{ini}_l(d), \mathrm{fin}_l(d)\,)$, que acaba al día siguiente si la salida
        no es posterior a la entrada (las noches)
  \item Su franja $\varphi(l) \in \Franjas = \{\text{mañana}, \text{tarde}, \text{noche},
        \text{partido}, \text{localizado}\}$ vienen definido en los datos
  \item Los minutos que computa para la jornada, $m(l) \in \mathbb{N}$ (no siempre coinciden con
        la duración del horario: un partido tiene una interrupción, y una guardia localizada de
        24~h computa 8h) el uso de minutos es para facilitar operaciones evitando punto flotante.
  \item Las personas que pide cada día, $p(l) \in \{1\}$ (por ahora solo he tratado turnos de 1).
\end{itemize}
La línea $l$ \emph{opera} el día $d$ si $\tau_{b(l)}(d) \in O(l)$, y entonces pide $p(l)$
personas ese día.

\subsection{La plantilla}

El conjunto de trabajadores $\Trab$ se divide en tres tipos:
\[
  \Trab = \Trab_{\mathrm{pat}} \;\cup\; \Trab_{\mathrm{fij}} \;\cup\; \Trab_{\mathrm{cor}},
\]
\begin{itemize}[nosep]
  \item \textbf{De patrón}: siguen una rotación pactada semana a semana (en la zona estudiada,
        las noches y las UVI)
  \item \textbf{Fijos}: trabajan en las líneas de su base o pertenecen al conjunto global que debe seguir una rotación en la que se garantice estabilidad pero a su vez equidad entre semanas de diferente tipo de turnos.
  \item \textbf{Correturnos}: cubren lo que queda, en cualquier base, con más flexibilidad
\end{itemize}
Cada trabajador $w$ tiene una base $b(w)$, unos días de vacaciones $V(w) \subseteq \Dias$ y una
jornada anual $H(w)$ en minutos (la del convenio, 1776~h, escalada si tiene reducción de
jornada).

\subsection{El cuadrante}

El resultado es una función parcial
\[
  P : \Trab \times \Dias \rightharpoonup \Lin \cup \Desc,
  \qquad \Desc = \{\mathrm{DS}, \mathrm{DO}\},
\]
donde $P(w,d) = l$ significa que $w$ hace la línea $l$ el día $d$, y los valores de $\Desc$ son
descansos con nombre: descanso semanal (DS) y descanso obligatorio tras una guardia (DO). Si $P(w,d)$ no está definida, ese día es un \emph{hueco}: no
tiene nada asignado y se puede rellenar. Un descanso no es un hueco: ocupa el día igual que un
turno.

Escribimos $t(w,d) = P(w,d)$ si $P(w,d) \in \Lin$, y $t(w,d) = \bot$ en otro caso (descansa o
está en hueco). A partir de $P$ se definen:
\begin{align*}
  \text{cobertura:}\quad & c(l,d) = \bigl|\{\, w \in \Trab : t(w,d) = l \,\}\bigr|
                           \;\le\; p(l), \text{l es línea,d día y }p(l) \text{ su demanda}\\
  \text{plazas sin cubrir:}\quad & h(l,d) = p(l) - c(l,d)
                           \quad \text{si $l$ opera el día $d$}, \\
  \text{minutos trabajados:}\quad & M(w) = \sum_{d \in \Dias,\; t(w,d) \neq \bot} m\bigl(t(w,d)\bigr).
\end{align*}

% =========================================================================== %
\section{Restricciones}
% =========================================================================== %

Las restricciones se dividen en dos familias. Las del \emph{convenio} son comunes a todas las
zonas que se rigen por él: el algoritmo las define una sola vez y solo cambian sus parámetros.
Las de \emph{organización} son propias de cada zona: cómo se reparten los descansos y los fines
de semana. Todas son \textbf{duras}: ningún objetivo puede comprarse a costa de incumplir una.

Hay una excepción deliberada: los \emph{patrones pactados} (noches y UVI) se dan por válidos
tal como vienen, porque están acordados con los trabajadores al menos esa es la casuística de Valladolid.
 Las restricciones se aplican a
todo lo que el algoritmo decide, no a los ciclos de patrón que se copian.

\subsection{Básicas}

Por construcción, $P$ asigna como mucho un valor a cada trabajador y día, así que nadie hace dos
turnos el mismo día. Además, para todo $w \in \Trab$, $d \in \Dias$ y $l \in \Lin$:
\begin{align*}
  & d \in V(w) \;\Rightarrow\; P(w,d) \text{ no está definida}
      && \text{(nadie trabaja ni descansa con nombre en vacaciones)}\\
  & t(w,d) = l \;\Rightarrow\; l \text{ opera el día } d
      && \\
  & c(l,d) \le p(l)
      && \text{(no se cubren más plazas de las que hay)}
\end{align*}
y quién puede hacer cada línea lo fija la sección~\ref{sec:quien}.

\subsection{Convenio}

El convenio fija cuatro parámetros: el descanso mínimo entre jornadas $\delta$ (12~h), el máximo
de días trabajados por semana $D_{\max}$ (6), el máximo de horas por semana $H_{\mathrm{sem}}$
(48~h) y la jornada anual (1776~h, de donde sale $H(w)$). Al menos así lo dicta el de CYL

\paragraph{C4 · Descanso entre jornadas.} Entre el fin de un turno y el inicio del siguiente, en
días consecutivos, median al menos $\delta$ horas:
\[
  t(w,d) = l,\;\; t(w,d+1) = l' \;\;\Rightarrow\;\;
  \mathrm{ini}_{l'}(d+1) - \mathrm{fin}_{l}(d) \;\ge\; \delta,
\]
salvo que $l$ o $l'$ sea una guardia localizada: es disponibilidad y no presencia, así que ni
exige descanso después ni lo consume antes. Solo hay que mirar días consecutivos: con un día
entero de por medio, el descanso se cumple solo.

\paragraph{C5 · Días por semana.} Para cada semana $s \in \Sem$:
\[
  \bigl|\{\, d \in s : t(w,d) \neq \bot \,\}\bigr| \;\le\; D_{\max}.
\]

\paragraph{C6 · Horas por semana.} Para cada semana $s \in \Sem$, en minutos:
\[
  \sum_{d \in s,\; t(w,d) \neq \bot} m\bigl(t(w,d)\bigr) \;\le\; 60\, H_{\mathrm{sem}}.
\]
C5 y C6 se cuentan por semana ISO, de lunes a domingo, y no en cualquier ventana de siete días.
Es la interpretación que se aplica en las zonas de este convenio, y admite, a sabiendas, rachas
por encima del tope a caballo entre dos semanas.

\paragraph{Jornada anual.} Nadie pasa de su jornada:
\[
  M(w) \;\le\; H(w).
\]
Es un techo duro. Quedarse lo más cerca posible de él es un \emph{objetivo}, no una restricción
(sección~\ref{sec:fijos}).

\subsection{Organización de la zona}

Para estas reglas hace falta saber qué días cuentan como \emph{libres}. Un día $d$ es libre para
$w$ si no trabaja ni está de vacaciones, y no es un festivo de lunes a viernes:
\[
  \mathrm{libre}(w,d) \iff t(w,d) = \bot,\;\; d \notin V(w),\;\;
  \neg\bigl(\tau_{b(w)}(d) = \mathrm{FES} \text{ y } d \text{ es de lunes a viernes}\bigr).
\]
Un festivo entre semana que no se trabaja es un descanso \emph{de más}, no el semanal; uno en
sábado o domingo sí cuenta como descanso semanal.

\paragraph{Descanso semanal.} En cada semana completa (sin vacaciones y dentro del año), al menos
dos días libres:
\[
  \bigl|\{\, d \in s : \mathrm{libre}(w,d) \,\}\bigr| \;\ge\; 2 .
\]
Al final se les pone nombre: son los DS del cuadrante. En algunas zonas estan fijado ciertos días, por ahora
no lo he implementado puesto que no he llegado.

\paragraph{Domingo con sábado.} Nadie trabaja un domingo sin trabajar el sábado de ese fin de
semana

\paragraph{Descanso de fin de semana.} Quien trabaja sábado y domingo tiene, esa misma semana,
dos días libres seguidos de lunes a viernes:
\[
  t(w,\mathrm{sab}_s) \neq \bot,\;\; t(w,\mathrm{dom}_s) \neq \bot \;\;\Rightarrow\;\;
  \exists\, d \in \{\mathrm{lun}_s, \dots, \mathrm{jue}_s\} :\;
  \mathrm{libre}(w,d) \,\wedge\, \mathrm{libre}(w,d+1).
\]

\subsection{Resumen}

\begin{center}
\begin{tabular}{@{}lll@{}}
  \toprule
  Regla & Ámbito & Qué limita \\
  \midrule
  Un turno al día, vacaciones, cobertura & básica & cada día \\
  C4 · descanso entre jornadas & convenio & cada par de días seguidos \\
  C5 · días por semana & convenio & cada semana ISO \\
  C6 · horas por semana & convenio & cada semana ISO \\
  Jornada anual & convenio & el año \\
  Descanso semanal & zona & cada semana completa \\
  Domingo con sábado & zona & cada fin de semana \\
  Descanso de fin de semana & zona & cada semana con finde trabajado \\
  \bottomrule
\end{tabular}
\end{center}

% =========================================================================== %
\section{Quién puede hacer cada línea}
\label{sec:quien}
% =========================================================================== %

No todo el mundo puede hacer cualquier línea. Para cada línea, los datos declaran quién la tiene
asignada y quién la cubre si falta. Esto generalmente viene por escrito o como anotación en algún excel
estoy contemplando realizar de alguna manera una interfaz que permita indicarlo porque recogerlo como fichero
me parece complejo. El resto se deduce de unas pocas reglas según el tipo de
trabajador. El resultado son tres \emph{grupos} que se consultan siempre en el mismo orden.

\subsection{Titulares, cubridores y pool}

Para cada línea $l$:
\begin{itemize}[nosep]
  \item $T(l) \subseteq \Trab$ son sus \textbf{titulares}: quienes la tienen asignada. Esos días
        la hacen ellos y no hacen otra. Los trabajadores de patrón son titulares de las líneas
        que recorre su patrón sin necesidad de declararlo.
  \item $C(l) \subseteq \Trab$ son sus \textbf{cubridores}: designados para cubrirla cuando
        falta un titular, en un orden de preferencia (principal y suplentes). Son el último
        recurso.
  \item $Q(l) \subseteq \Trab$ es su \textbf{pool}: quién más puede hacerla por la regla de su
        tipo. Con los fijos de su base que no son titulares de nada y los correturnos que la
        admiten,
        \[
          Q_{\mathrm{fij}}(l) = \{\, w \in \Trab_{\mathrm{fij}} \setminus \mathcal{T} : b(w) = b(l) \,\},
          \qquad
          Q_{\mathrm{cor}}(l) = \{\, w \in \Trab_{\mathrm{cor}} : \mathrm{admite}(w,l) \,\},
        \]
        el pool es
        \[
          Q(l) =
          \begin{cases}
            \varnothing & \text{si } C(l) \neq \varnothing,\\[2pt]
            Q_{\mathrm{cor}}(l) & \text{si } C(l) = \varnothing \text{ y } T(l) \neq \varnothing,\\[2pt]
            Q_{\mathrm{fij}}(l) \cup Q_{\mathrm{cor}}(l) & \text{en otro caso,}
          \end{cases}
        \]
        donde $\mathcal{T} = \bigcup_{l'} T(l')$ son los titulares de alguna línea (quien tiene
        líneas asignadas solo hace las suyas), y un correturno no hace noches ni guardias
        localizadas de una base que no es la suya pero si mañanas,tardes y partido:
        \[
          \mathrm{admite}(w,l) \iff \neg\bigl(\varphi(l) \in \{\text{noche}, \text{localizado}\}
          \;\wedge\; b(l) \neq b(w)\bigr).
        \]
\end{itemize}
En palabras: una línea con cubridores es exclusiva de sus titulares y sus cubridores; una línea
con titulares y sin cubridores la hacen sus titulares y, cuando faltan, los correturnos; y una
línea sin titulares es de todo el pool de su base.

\begin{center}
\begin{tabular}{@{}llll@{}}
  \toprule
  Línea & Titulares & Cubridores & Pool \\
  \midrule
  sin titulares ni cubridores & --- & --- & fijos de su base y correturnos \\
  con titulares & los declarados & --- & solo correturnos \\
  con titulares y cubridores & los declarados & los designados & nadie \\
  de patrón & los del patrón & los designados & nadie \\
  \bottomrule
\end{tabular}
\end{center}

\subsection{Conciliaciones}

Algunos trabajadores tienen \emph{restricciones} propias: conciliaciones o días que no trabajan.
Cada restricción $r$ de $w$ es una terna $(T_r, L_r, J_r)$: los tipos de día a los que se aplica
$T_r \subseteq \Tipos$, las líneas que esos días puede hacer $L_r \subseteq \Lin$ (o todas), y la
ventana horaria $J_r$ dentro de la que tiene que caer el turno (o ninguna). Por ejemplo,
\emph{``de lunes a viernes, solo de 8:00 a 16:00''} o \emph{``sábados, domingos y festivos,
ninguna línea''}. Que $w$ \emph{pueda} hacer $l$ el día $d$ es:
\[
  \mathrm{puede}(w,l,d) \iff d \notin V(w) \;\wedge\;
  \forall r : \tau_{b(l)}(d) \in T_r \Rightarrow
  \bigl( l \in L_r \;\wedge\; [\,\mathrm{ini}_l(d), \mathrm{fin}_l(d)\,) \subseteq J_r \bigr).
\]
Las restricciones solo quitan opciones; nunca dan acceso a una línea.

\subsection{Los grupos de un día}

Juntando todo, quién puede hacer la línea $l$ el día $d$ son tres grupos, en orden. Si $l$ opera
el día $d$,
\[
  \mathrm{grupos}(l,d) = \bigl( T_d(l),\; C_d(l),\; Q_d(l) \bigr),
  \qquad
  X_d(l) = \{\, w \in X(l) : \mathrm{puede}(w,l,d) \,\},
\]
y tres conjuntos vacíos si no. Los pasos del algoritmo consultan los
grupos en ese orden: se pasa al siguiente solo si en el anterior no hay nadie a quien poner
\emph{legalmente} con lo que ya tiene en el cuadrante. Así, un titular que ese día no puede
porque incumpliría el descanso de C4 no bloquea la línea: la cubre un cubridor o alguien del
pool.

Esta función responde a \emph{quién puede}, no a \emph{si es legal} (eso son las
restricciones de la sección anterior, que dependen del resto del cuadrante) ni a \emph{a quién
se prefiere} dentro de cada grupo (eso lo decide cada paso).

% =========================================================================== %
\section{Arquitectura y orden de los pasos}
% =========================================================================== %

\subsection{Tres capas}

El generador separa tres cosas que cambian a ritmos distintos:
\begin{itemize}[nosep]
  \item \textbf{El convenio}: los límites legales de la sección 3 (C4, C5, C6, jornada). Es
        común a todas las zonas de una comunidad y cambia cada pocos años. En el código es un
        objeto con sus parámetros, no un dato que se pueda escribir mal.
  \item \textbf{La zona}: cómo se organiza el trabajo en ella. Sus reglas de reparto (descanso
        semanal, fines de semana) y su \emph{receta}: qué pasos se
        dan y en qué orden. Cada zona es un problema independiente, con su plantilla y sus
        líneas.
  \item \textbf{Los datos}: el plan funcional, la plantilla y los festivos de un año.
\end{itemize}
Con esta separación, una zona nueva no obliga a rehacer lo anterior: se añade su receta,
reutilizando los pasos que comparta con otras, y las restricciones del convenio y la salida
siguen siendo las mismas.

\subsection{Los pasos}

El algoritmo construye el cuadrante $P$ por pasos. Cada paso recibe el cuadrante que deja el
anterior, lo trata como constante (lo ya asignado ocupa el día y gasta jornada) y lo completa.
La figura~\ref{fig:pasos} muestra el orden en la zona estudiada.

\begin{figure}[ht]
\centering
\begin{tikzpicture}[
    node distance=5mm and 8mm,
    paso/.style={draw, rounded corners=2pt, align=center, text width=6.6cm,
                 minimum height=9mm, font=\small},
    muestra/.style={draw, rounded corners=2pt, minimum width=8mm, minimum height=4mm},
    regla/.style={paso, fill=gray!12},
    optim/.style={paso, fill=blue!12, draw=blue!50!black},
    flecha/.style={-{Stealth[length=2.5mm]}, thick},
    nota/.style={font=\footnotesize, align=left, text width=4.6cm}
  ]
  \node[regla] (carga)  {Carga de datos\\{\footnotesize plan funcional, plantilla, festivos}};
  \node[regla, below=of carga] (esq) {Esqueleto\\{\footnotesize patrones pactados; vacaciones como ausencia}};
  \node[regla, below=of esq] (lib) {Libranzas\\{\footnotesize ausencias y exceso de horas de los patrones}};
  \node[regla, below=of lib] (fin) {Fines de semana y festivos\\{\footnotesize reparto greedy por turnos}};
  \node[optim, below=of fin] (fij) {Fijos, de lunes a viernes\\{\footnotesize CP-SAT anual, una base cada vez}};
  \node[optim, below=of fij] (cor) {Correturnos\\{\footnotesize CP-SAT anual sobre lo que queda}};
  \node[regla, below=of cor] (des) {Descanso semanal\\{\footnotesize DS en su día}};
  \node[regla, below=of des] (sal) {Comprobación del convenio y Excel};

  \foreach \a/\b in {carga/esq, esq/lib, lib/fin, fin/fij, fij/cor, cor/des, des/sal}
    \draw[flecha] (\a) -- (\b);

  \begin{scope}[on background layer]
    \node[draw=black!60, dashed, rounded corners=4pt, inner sep=3mm,
          fit=(lib)(fin)(fij)(cor)(des)] (zona) {};
  \end{scope}

  \node[nota, right=8mm of esq.east |- carga] {\textbf{Común} a todas las zonas};
  \node[nota, right=8mm of zona.east] {\textbf{Receta de la zona}: qué pasos y en qué orden;
        otra zona puede tener otros};
  \node[nota, right=8mm of esq.east |- sal] {\textbf{Común}: el convenio y la salida no
        dependen de la zona};

  \node[muestra, fill=gray!12, below=9mm of sal.south west, anchor=north west] (lr) {};
  \node[right=1mm of lr, font=\footnotesize] (tr) {regla determinista};
  \node[muestra, fill=blue!12, draw=blue!50!black, right=8mm of tr] (lo) {};
  \node[right=1mm of lo, font=\footnotesize] {modelo de optimización};
\end{tikzpicture}
\caption{Los pasos del algoritmo en la zona estudiada. Los del recuadro forman la receta de la
zona; los demás son comunes.}
\label{fig:pasos}
\end{figure}

El orden no es arbitrario: cada paso coloca antes lo que está más atado, para que los
siguientes, que tienen más libertad, se adapten a ello.
\begin{enumerate}[nosep]
  \item El \textbf{esqueleto} pinta lo que ya viene decidido: los ciclos de los patrones
        pactados.
  \item Las \textbf{libranzas} resuelven lo que esos ciclos no pueden cumplir por sí solos: quién
        hace el ciclo de un titular de vacaciones y cómo devuelve un patrón las horas que pasa de
        su jornada.
  \item Los \textbf{fines de semana y festivos} se reparten antes que el lunes a viernes porque
        son la parte más escasa y la que más pesa en la equidad, y porque determinan qué semanas
        necesitan libres entre semana.
  \item Los \textbf{fijos} se resuelven antes que los correturnos porque necesitan estabilidad
        (semanas enteras de la misma franja) y llegar a su jornada; los correturnos, que son la
        parte flexible de la plantilla, recogen lo que queda.
  \item Los \textbf{descansos semanales} se escriben al final, cuando ya se sabe qué días quedan
        libres. Los modelos anteriores ya garantizan que haya sitio para ellos.
\end{enumerate}

% =========================================================================== %
\section{Pasos deterministas}
% =========================================================================== %

Los tres primeros pasos no optimizan nada: aplican reglas fijas, en un orden fijo, y con los mismos
datos dan siempre el mismo cuadrante.

\subsection{Esqueleto: los patrones pactados}

Un patrón $\pi$ es una matriz $A_\pi$ de $R_\pi$ filas por 7 columnas: cada fila es una semana
de lunes a domingo, y cada celda es una línea, un descanso de $\Desc$ o nada. Todos los
trabajadores del patrón recorren las mismas filas, una por semana, cada uno desde la suya,
$f_0(w)$, que es la fila que le toca en la primera semana del año y da continuidad de un año al
siguiente.

Si $k(d) = \lfloor (d - d_0)/7 \rfloor$ es el número de semana del día $d$ contado desde el lunes
$d_0$ de la semana del 1 de enero, y $j(d) \in \{0,\dots,6\}$ su día de la semana, lo que el
patrón prescribe a $w$ ese día es
\[
  \mathrm{pres}(w,d) = A_\pi\bigl[\,(f_0(w) + k(d)) \bmod R_\pi,\; j(d)\,\bigr].
\]
El esqueleto es
\[
  P(w,d) = \mathrm{pres}(w,d)
  \qquad \text{si } d \notin V(w) \text{ y } \bigl(\mathrm{pres}(w,d) \in \Desc
  \text{ o la línea opera el día } d\bigr),
\]
y nada en otro caso. Las vacaciones no se pintan: son ausencia. Tampoco se comprueba ninguna
restricción, porque los patrones se dan por válidos (sección~3).

\subsection{Libranzas: ausencias y exceso de horas}

Los patrones dejan dos problemas que el esqueleto no resuelve. En la zona estudiada todas las
líneas de patrón tienen cubridores, así que los dos se resuelven del mismo modo: un cubridor
\emph{hereda} días del titular, haciendo lo que el patrón le prescribe a este, turnos y
descansos:
\[
  P(c,d) \leftarrow \mathrm{pres}(w,d) \qquad \text{para los días } d \text{ heredados}.
\]
Lo que se hereda nunca se comprueba: si el ciclo era válido para el titular, lo es para quien lo
sustituye.

\paragraph{Vacaciones del titular.} Los días de $V(w)$, el patrón de $w$ queda sin nadie. Se
hereda a un cubridor el periodo entero. Se elige el primero, en el orden de preferencia, que
esté \emph{intacto} esos días (disponible y sin que nadie haya tocado lo suyo) y que no venga de
cubrir algo a menos de 15 días; si todos vienen de cubrir algo cerca, el que lo tenga más lejos.
Esa separación es la que reparte las coberturas entre el principal y los suplentes: sin ella, el
principal encadena una quincena tras otra y el suplente no entra nunca. Si nadie puede el
periodo entero, se reparte en tramos.

\paragraph{Exceso de horas.} Un patrón suele prescribir más horas de las de la jornada, porque
no está hecho a la medida del año: $\sum_{d} m(\mathrm{pres}(w,d))$ puede superar $H(w)$. El
exceso
\[
  e(w) = M(w) - H(w) > 0
\]
se \emph{cede}: el titular deja de hacer algunos días, que hereda un cubridor. La unidad de cesión
es un \emph{bloque}: una racha máxima de días de trabajo del patrón seguida de los descansos que
la compensan. Ceder bloques enteros respeta la forma del ciclo, y lo que hereda el cubridor es un
trozo con sentido (un turno de noches con sus libres), no días sueltos.

Mientras $e(w) > 0$, se cede un bloque más. Entre los bloques que siguen intactos y para los que
hay un cubridor libre, se elige por dos criterios, en orden:
\begin{enumerate}[nosep]
  \item que el bloque caiga en días en que la base tenga \emph{holgura}, porque el cubridor que
        se va deja un hueco que alguien tendrá que cubrir. La holgura de un día estima cuánto
        le sobra a la plantilla de la base:
        \[
          \mathrm{holg}(d) = \rho \cdot \bigl|\{\, w' \text{ libres el día } d \,\}\bigr|
                             \;-\; \mathrm{dem}(d),
          \qquad \rho = \frac{1776}{8 \cdot 365} \approx 0{,}6,
        \]
        donde $\rho$ es la fracción de días que trabaja de media una persona y
        $\mathrm{dem}(d)$ las plazas de la base que ese día puede cubrir esa plantilla. Se prefiere
        un bloque con holgura media no negativa;
  \item a igualdad, el que quede más lejos de lo que el cubridor ya tiene comprometido.
\end{enumerate}

\paragraph{Recorte del último bloque.} Ceder siempre bloques enteros puede dejar al titular muy
por debajo de su jornada: un bloque de noches son unas 60 horas. Por eso, si la racha entera
quita más de lo que sobra, solo se cede su final. Si la racha son los días $d_1 < \dots < d_n$
con minutos $m_1, \dots, m_n$ y $\sum_i m_i > e(w)$, se ceden los días $d_j, \dots, d_n$ con
\[
  j = \max\Bigl\{\, j : \sum_{i=j}^{n} m_i \ge e(w) \,\Bigr\},
\]
es decir, las últimas jornadas justas para no pasar de la jornada. El titular acorta su racha y
conserva sus descansos, y el cubridor hace el final. Así cada titular acaba a menos de un turno
por debajo de $H(w)$.

\subsection{Fines de semana y festivos}

Sábados, domingos y festivos se reparten antes que el lunes a viernes, entre fijos y
correturnos, con una regla \emph{greedy}: se recorre el año en orden y cada plaza se da al
candidato que menos lleva de esa clase. Es un turno rotatorio: quien acaba de trabajar un
domingo pasa al final de la cola de los domingos, y no le vuelve a tocar hasta que los demás le
alcanzan. Si todos estuvieran siempre disponibles, sería un turno exacto y los recuentos de un
mismo grupo diferirían como mucho en uno, las vacaciones, las conciliaciones y la legalidad lo
alteran, pero la regla siempre corrige hacia la igualdad. Además, como la cola obliga a esperar
el turno, los fines de semana de cada persona quedan espaciados por el año sin medirlo
explícitamente.

Se hacen tres pasadas, una por clase de día $\kappa$, en el orden
$\mathrm{FES} \to \mathrm{DOM} \to \mathrm{SAB}$. Un sábado festivo es de clase $\mathrm{FES}$ y
así, cuando llega su domingo, ya tiene gente.

\paragraph{Candidatos.} Para la plaza de la línea $l$ el día $d$, los candidatos son los del
primer grupo de $\mathrm{grupos}(l,d)$ (sección~\ref{sec:quien}) en el que hay alguien a quien se
le puede poner $l$ ese día sin incumplir ninguna restricción con lo que ya tiene en el
cuadrante, ni pasar de su jornada.

\paragraph{El sábado va con su domingo.} La regla del domingo con sábado no se comprueba
después: se cumple por construcción. Si a quien coge un domingo no le toca trabajar ese sábado,
se le asigna a la vez una plaza pendiente de ese sábado, buscando primero la misma línea. Si no
hay ninguna con la que el domingo sea legal, no es candidato para ese domingo. Esos sábados ya
cuentan cuando empieza la pasada de los sábados.

\paragraph{Qué plaza y a quién.} Cada día se cubre primero la plaza con \emph{menos candidatos},
la misma idea que la heurística de la variable más restringida en satisfacción de restricciones:
no gastar en una plaza fácil a alguien que era el único posible para otra. Entre los candidatos
se elige el mínimo, en orden lexicográfico, de
\[
  \mathrm{clave}(w) = \bigl(\, \neg\,\mathrm{misma}(w),\;\; n_\kappa(w),\;\; -(d - u(w)),\;\;
  n_{\mathrm{FES}}(w) + n_{\mathrm{DOM}}(w) + n_{\mathrm{SAB}}(w) \,\bigr),
\]
donde
\begin{itemize}[nosep]
  \item $\mathrm{misma}(w)$, solo los domingos: si con esta asignación $w$ hace la misma línea el
        sábado y el domingo. Así el fin de semana de una guardia localizada es de una sola
        persona, sin tratarlo como caso aparte
  \item $n_\kappa(w)$: cuántos días de la clase $\kappa$ lleva ya, contando los que trae el
        cuadrante de los pasos anteriores (los ciclos heredados);
  \item $u(w)$: el último fin de semana que trabajó, para preferir a quien lleva más tiempo sin
        trabajar uno;
  \item el total de fines de semana y festivos, como último desempate.
\end{itemize}

\begin{algorithm}[ht]
\caption{Reparto de sábados, domingos y festivos}
\label{alg:findes}
\begin{algorithmic}[1]
  \Require el cuadrante $P$ tras las libranzas
  \For{$\kappa \in (\mathrm{FES}, \mathrm{DOM}, \mathrm{SAB})$}
    \For{cada día $d$ de clase $\kappa$, en orden}
      \While{quedan plazas pendientes de clase $\kappa$ el día $d$}
        \State $\Omega(l) \gets$ candidatos de cada plaza pendiente $l$, cada uno con su sábado
               si $\kappa = \mathrm{DOM}$
        \State $l^\ast \gets \arg\min_l |\Omega(l)|$
          \Comment{la plaza más restringida}
        \If{$\Omega(l^\ast) = \varnothing$}
          \State marcar $l^\ast$ como sin cubrir
        \Else
          \State $w \gets \arg\min_{w \in \Omega(l^\ast)} \mathrm{clave}(w)$
          \State $P(w,d) \gets l^\ast$, y su plaza de sábado si la lleva
          \State actualizar $n_\kappa(w)$ y $u(w)$
        \EndIf
      \EndWhile
    \EndFor
  \EndFor
\end{algorithmic}
\end{algorithm}

Las líneas de patrón no entran en este paso: ya las cubren el esqueleto y las libranzas. Lo que
el reparto no consigue cubrir no se reintenta: queda como plaza sin cubrir.

% =========================================================================== %
\section{El modelo de los fijos}
\label{sec:fijos}
% =========================================================================== %

Con los fines de semana y festivos repartidos, queda el lunes a viernes: el grueso del año y la
parte que más condiciona la vida de la plantilla fija, que quiere semanas estables y un reparto
justo de las semanas duras. Es el primer paso que necesita ver el año entero a la vez, porque la
jornada, las cuotas y la rotación son restricciones anuales: una decisión de
marzo cambia lo que se puede hacer en noviembre. Se plantea como un problema de programación con
restricciones sobre variables booleanas y se resuelve con CP-SAT.

\subsection{La idea, en palabras}

Antes de las fórmulas, conviene ver qué pregunta responde el modelo. Para cada fijo y cada día
laborable del año, decide \textbf{si trabaja y en qué tipo de trabajo}: por ejemplo, ``mañana de
8 horas en Valladolid''. No decide todavía la línea concreta, porque para el resto del año da
igual cuál de las mañanas de 8 horas haga. Eso se resuelve después, semana a semana.

Sobre esa decisión se imponen las reglas de la sección 3, y una más que es propia de este
modelo: \textbf{cada semana, un solo tipo de trabajo}. Así la semana sale estable por
construcción: quien hace mañanas un lunes las hace toda la semana.

Lo que ya dejaron los pasos anteriores (los sábados y domingos repartidos, los ciclos heredados
de un patrón, las vacaciones) no se toca: entra en el modelo como dato. Si alguien trabaja el
sábado, ese sábado cuenta para sus días de la semana, sus horas y su descanso.

La figura~\ref{fig:semana} muestra una semana de un fijo tal como la ve el modelo.

\begin{figure}[ht]
\centering
\begin{tikzpicture}[
    x=1.05cm, y=0.62cm,
    celda/.style={draw=black!40, minimum width=1.05cm, minimum height=0.62cm, anchor=center},
    si/.style={celda, fill=blue!25},
    fijo/.style={celda, fill=gray!30},
    nada/.style={celda, fill=white},
    fuera/.style={celda, pattern=north east lines, pattern color=black!25},
    rot/.style={font=\footnotesize, anchor=east}
  ]
  % cabecera de días
  \foreach \d [count=\i] in {L,M,X,J,V,S,D}
    \node[font=\small\bfseries] at (\i,1) {\d};
  % fila: lo que ya viene dado
  \node[rot] at (0.4,0) {ya decidido};
  \foreach \i in {1,...,5} \node[nada] at (\i,0) {};
  \node[fijo, font=\scriptsize] at (6,0) {finde};
  \node[nada] at (7,0) {};
  % filas de casillas
  \node[rot] at (0.4,-1.3) {mañana · 8 h};
  \node[rot] at (0.4,-2.3) {tarde · 8 h};
  \node[rot] at (0.4,-3.3) {partido · 10 h};
  \foreach \i/\v in {1/1,2/1,3/0,4/1,5/1} {
    \ifnum\v=1 \node[si, font=\scriptsize] at (\i,-1.3) {1}; \else \node[nada, font=\scriptsize] at (\i,-1.3) {0}; \fi
    \node[nada, font=\scriptsize] at (\i,-2.3) {0};
    \node[nada, font=\scriptsize] at (\i,-3.3) {0};
  }
  \foreach \i in {6,7} \foreach \r in {-1.3,-2.3,-3.3} \node[fuera] at (\i,\r) {};
  % columna de zona
  \node[font=\small\bfseries] at (8.6,1) {zona};
  \node[si, font=\scriptsize] at (8.6,-1.3) {1};
  \node[nada, font=\scriptsize] at (8.6,-2.3) {0};
  \node[nada, font=\scriptsize] at (8.6,-3.3) {0};
  % fila de trabaja
  \node[rot] at (0.4,-4.6) {trabaja $x_{w,d}$};
  \foreach \i/\v in {1/1,2/1,3/0,4/1,5/1,6/1,7/0}
    \node[nada, font=\scriptsize] at (\i,-4.6) {\v};
  % llaves y notas
  \node[font=\footnotesize, align=left, anchor=west] at (9.3,-1.3) {una sola zona\\en la semana};
  \node[font=\footnotesize, align=left, anchor=west] at (9.3,-4.6) {5 días ($\le 6$),\\libres: X y D};
\end{tikzpicture}
\caption{Una semana de un fijo vista por el modelo. Arriba, lo que dejaron los pasos anteriores
(el sábado de su fin de semana). En medio, las variables $y$: un 1 en una casilla y día es que
ese día trabaja en ese tipo de trabajo. Los fines de semana no tienen variables (rayado): ya
están decididos. A la derecha, la zona de la semana: solo una puede valer 1, y todas las
casillas marcadas tienen que ser de ella. Abajo, lo que trabaja cada día, que es lo que miran
C5, C6, la jornada y el descanso.}
\label{fig:semana}
\end{figure}

Entre todas las formas de rellenar el año que cumplen las reglas, el modelo elige por orden:
primero la que cubre más plazas; entre esas, la que deja a cada fijo más cerca de su jornada;
y entre esas, la que reparte mejor las semanas duras. El resto de la sección precisa cada una de
estas piezas.

\subsection{Descomposición}

Un fijo solo hace líneas de su base, así que \textbf{cada base es un problema independiente}: se
resuelve un modelo por base, y uno pequeño se resuelve mejor y antes que uno grande. Dentro de
cada base el problema se divide en dos fases:
\begin{enumerate}[nosep]
  \item un \textbf{modelo anual} que decide, para cada fijo y cada día laborable, \emph{en qué
        tipo de trabajo} está, sin elegir la línea concreta;
  \item una \textbf{asignación directa} que, semana a semana, elige la línea concreta dentro de
        ese tipo.
\end{enumerate}

\subsection{Casillas y zonas}

Ese ``tipo de trabajo'' es lo que llamamos \emph{casilla}. Muchas líneas son intercambiables para
todas las restricciones del año: dos mañanas de la misma base que computan lo mismo gastan la
misma jornada, cuentan igual en C5 y C6 y pesan lo mismo en las cuotas. Decidir entre ellas
dentro del modelo multiplicaría las variables y, sobre todo, las soluciones simétricas, sin
mejorar nada.

Para agrupar líneas se usan tres datos. La \emph{bolsa} $\beta(l)$ de una línea es el conjunto
de líneas que rota la misma gente: si $l$ no tiene titulares, las abiertas de su base; si los
tiene, las de ese mismo grupo de titulares. Con ella,
\[
  \mathrm{cas}(l) = \bigl(\beta(l),\; \varphi(l),\; m(l)\bigr) \in K,
  \qquad
  \zeta(l) = \bigl(\beta(l),\; \varphi(l)\bigr),
\]
son su \emph{casilla} (bolsa, franja y minutos) y su \emph{zona} (bolsa y franja). La casilla es
lo que se decide cada día; la zona es lo que tiene que ser igual toda la semana. Los minutos
separan las líneas que no computan lo mismo, para que la jornada y C6 salgan exactas; la zona
los ignora, porque hacer una mañana de 7 horas el lunes y una de 9 el martes sigue siendo una
semana de mañanas (figura~\ref{fig:casillas}).

\begin{figure}[ht]
\centering
\begin{tikzpicture}[
    node distance=3mm and 14mm,
    lin/.style={draw, rounded corners=2pt, font=\small, minimum width=2.6cm, minimum height=6mm},
    cas/.style={draw, rounded corners=2pt, font=\small, minimum width=3.1cm, minimum height=6mm, fill=blue!10},
    zon/.style={draw, rounded corners=2pt, font=\small, minimum width=2.6cm, minimum height=6mm, fill=blue!25},
    flecha/.style={-{Stealth[length=2mm]}, black!60}
  ]
  \node[font=\small\bfseries] at (0,0.9) {Líneas};
  \node[font=\small\bfseries] at (4.9,0.9) {Casillas};
  \node[font=\small\bfseries] at (9.0,0.9) {Zonas};
  \node[lin] (l1) at (0,0)    {A · mañana · 7 h};
  \node[lin] (l2) at (0,-0.9) {B · mañana · 8 h};
  \node[lin] (l3) at (0,-1.8) {C · mañana · 8 h};
  \node[lin] (l4) at (0,-2.7) {D · mañana · 9 h};
  \node[lin] (l5) at (0,-3.8) {E · tarde · 8 h};
  \node[lin] (l6) at (0,-4.7) {F · tarde · 8 h};
  \node[cas] (c1) at (4.9,0)     {mañana · 7 h};
  \node[cas] (c2) at (4.9,-1.35) {mañana · 8 h};
  \node[cas] (c3) at (4.9,-2.7)  {mañana · 9 h};
  \node[cas] (c4) at (4.9,-4.25) {tarde · 8 h};
  \node[zon] (z1) at (9.0,-1.35) {mañana};
  \node[zon] (z2) at (9.0,-4.25) {tarde};
  \foreach \a/\b in {l1/c1, l2/c2, l3/c2, l4/c3, l5/c4, l6/c4} \draw[flecha] (\a.east) -- (\b.west);
  \foreach \a/\b in {c1/z1, c2/z1, c3/z1, c4/z2} \draw[flecha] (\a.east) -- (\b.west);

  \node[font=\footnotesize, align=center] at (0,-5.7)   {se elige en la\\segunda fase};
  \node[font=\footnotesize, align=center] at (4.9,-5.7) {se decide\\cada día ($y$)};
  \node[font=\footnotesize, align=center] at (9.0,-5.7) {una por\\semana ($z$)};
\end{tikzpicture}
\caption{Un ejemplo con seis líneas ficticias de una misma bolsa. Las dos mañanas de 8 horas son
la misma casilla; las cuatro mañanas, la misma zona.}
\label{fig:casillas}
\end{figure}

Para escribir el modelo hacen falta las plazas pendientes y, para cada fijo, sus opciones. Sea
$n(l,d)$ el número de plazas de la línea $l$ que el día laborable $d$ siguen sin cubrir tras los
pasos anteriores, sin contar las de patrón. Las plazas de una casilla son
$N(k,d) = \sum_{\mathrm{cas}(l)=k} n(l,d)$. Las \emph{opciones} de un fijo $w$ el día $d$ en la
casilla $k$ son las líneas de esa casilla con plaza libre que puede hacer:
\[
  \Lambda(w,d,k) = \bigl\{\, l : \mathrm{cas}(l) = k,\;\; n(l,d) > 0,\;\;
    w \in T_d(l) \cup C_d(l) \cup Q_d(l) \,\bigr\},
\]
siempre que $P(w,d)$ esté vacío y $l$ cumpla C4 con lo que ya tiene el día anterior y el
siguiente.

\subsection{Variables}

Son las dos capas de la figura~\ref{fig:semana}. Para cada $(w,d,k)$ con
$\Lambda(w,d,k) \neq \varnothing$, una booleana
\[
  y_{w,d,k} = 1 \iff w \text{ trabaja el día } d \text{ en una línea de la casilla } k,
\]
y para cada fijo, semana y zona posible, una booleana $z_{w,s,\zeta}$: si esa semana trabaja en
la zona $\zeta$. Con ellas, que $w$ trabaje el día $d$ es
\[
  x_{w,d} =
  \begin{cases}
    1 & \text{si } P(w,d) \in \Lin \text{ ya estaba decidido},\\
    0 & \text{si } P(w,d) \in \Desc,\\
    \sum_k y_{w,d,k} & \text{si } P(w,d) \text{ está vacío,}
  \end{cases}
\]
es decir, lo que dejaron los pasos anteriores entra como constante. En la base de Valladolid
salen unas 20.000 variables $y$ y 4.400 variables $z$ para 51 fijos.

\subsection{Restricciones}

Son las de la sección 3, escritas sobre estas variables, más la de una zona por semana. Para
cada fijo $w$:

\paragraph{Un turno al día y cobertura.} Cada día marca como mucho una casilla, y en cada casilla
no entran más personas que plazas libres tiene:
\[
  \sum_k y_{w,d,k} \le 1,
  \qquad
  \sum_w y_{w,d,k} \le N(k,d).
\]
Además, quienes solo pueden hacer una línea $l$ de la casilla no pueden ser más que las plazas de
esa línea: sin esto el recuento por casilla cuadraría, pero la segunda fase no tendría dónde
ponerlos.

\paragraph{Una zona por semana.} Una casilla solo se puede marcar si su zona es la de la semana,
la zona solo vale 1 si de verdad se trabaja en ella, y cada semana tiene como mucho una zona:
\[
  y_{w,d,k} \le z_{w,s(d),\zeta(k)}, \qquad
  z_{w,s,\zeta} \le \sum_{d \in s,\; \zeta(k) = \zeta} y_{w,d,k}, \qquad
  \sum_\zeta z_{w,s,\zeta} \le 1.
\]
La segunda desigualdad evita que los niveles de cuotas marquen semanas vacías para cuadrar. Hay
una excepción: los días en que $w$ solo puede entrar como \emph{cubridor} de una línea no forman
zona, porque cubrir no es parte de ningún reparto. Sin esa excepción, un cubridor que tiene que
cubrir un día suelto perdería el resto de su semana.

\paragraph{C4 entre días seguidos.} Dos casillas en días consecutivos solo se prohíben si ninguna
de sus líneas descansa lo suficiente entre una y otra:
\[
  y_{w,d,k} + y_{w,d+1,k'} \le 1
  \quad \text{si ningún par } (l,l') \in \Lambda(w,d,k) \times \Lambda(w,d+1,k')
  \text{ cumple C4.}
\]
La restricción es \emph{optimista}: basta con que exista un par de líneas compatible. Lo exacto
lo comprueba la segunda fase al elegir la línea.

\paragraph{Semana: C5, C6 y descansos.} Se cuentan los días y los minutos de cada semana,
incluidos los que ya venían decididos:
\[
  \sum_{d \in s} x_{w,d} \le D_{\max},
  \qquad
  \sum_{d \in s} \Bigl( \sum_k m(k)\, y_{w,d,k} + m_0(w,d) \Bigr) \le 60\, H_{\mathrm{sem}},
\]
donde $m(k)$ son los minutos de la casilla y $m_0(w,d)$ los de lo que ya estaba en el
cuadrante. El descanso semanal y el de fin de semana se escriben con los indicadores
$1 - x_{w,d}$ de los días que cuentan como libres (sección 3); el par de libres seguidos usa una
booleana auxiliar por cada par de días consecutivos.

\paragraph{Jornada.} Nadie pasa de su jornada, y lo que le falta para llegar es su
\emph{déficit} $\delta_w$. Con $M_0(w)$ los minutos que ya trae del cuadrante,
\[
  M_0(w) + \sum_{d,k} m(k)\, y_{w,d,k} \;\le\; H(w),
  \qquad
  \delta_w = H(w) - M_0(w) - \sum_{d,k} m(k)\, y_{w,d,k} \;\ge\; 0.
\]

\subsection{Cuotas de semanas duras}

Las franjas \emph{duras} son la noche, el partido y la tarde. La idea es sencilla: si un tercio
de las plazas del año de una bolsa son de tarde, a cada fijo de esa bolsa le toca, más o menos,
un tercio de sus semanas de tarde. Eso es su \emph{cuota}.

Formalmente, cada fijo que puede elegir entre varias franjas de su bolsa tiene una cuota de cada
franja dura. Si $\pi(\beta,\varphi)$ es la proporción de plazas de la franja $\varphi$ en la
bolsa $\beta$ y $a_w$ el número de semanas en que $w$ está disponible,
\[
  q_{w,\varphi} = \bigl\lfloor \pi(\beta(w), \varphi) \cdot a_w \bigr\rfloor,
  \qquad
  Z_{w,\varphi} = \sum_s z_{w,s,(\beta(w),\varphi)},
\]
y se mide lo que se aleja de ella, por arriba o por abajo, con
$\mathrm{dev}_{w,\varphi} \ge |Z_{w,\varphi} - q_{w,\varphi}|$ (dos desigualdades lineales).

\paragraph{Ejemplo.} En la bolsa abierta de Valladolid, el 36\,\% de las plazas de lunes a
viernes son de tarde. Un fijo disponible 48 semanas tiene una cuota de
$\lfloor 0{,}36 \cdot 48 \rfloor = \lfloor 17{,}3 \rfloor = 17$ semanas de tarde.

El redondeo hacia abajo es deliberado: las semanas que sobran se quedan para los correturnos, en
la misma proporción, y así no se les cargan solo las duras. Si una franja no la puede hacer
ningún correturno (las noches de Medina), no hay a quién dejar las semanas que sobran. Entonces
la cuota es la parte proporcional de las semanas que hagan falta en total, y con
$A = \sum_v a_v$ se pide
\[
  \mathrm{dev}_{w,\varphi} \;\ge\; \Bigl|\, A \cdot Z_{w,\varphi} - a_w \sum_v Z_{v,\varphi} \,\Bigr|,
\]
que es lo que se aleja de su parte, multiplicado por $A$ para no usar números racionales.

Quien solo puede una franja (titular de una única línea o con una conciliación) no tiene cuota:
su semana ya viene dada.

\subsection{Rotación}

Cumplir la cuota no basta: siete semanas de tarde seguidas cumplen la cuota y son un mal
cuadrante. La plantilla espera una rotación reconocible, y se impone como restricción dura. La
idea es repartir las semanas duras de cada persona a lo largo del año dejando entre ellas, como
poco, la distancia que les tocaría si estuvieran espaciadas por igual.

Para cada fijo con cuota $q > 0$ en una franja dura $\varphi$, se fija la \emph{ventana}
\[
  g_{w,\varphi} = \max\bigl(2,\; \lfloor a_w / q_{w,\varphi} \rfloor\bigr)
\]
y se pide como mucho una semana de esa franja en cada $g$ semanas seguidas:
\[
  \sum_{j=i}^{i+g-1} z_{w,s_j,(\beta(w),\varphi)} \;\le\; 1
  \qquad \text{para todo } i.
\]
Con 7 semanas de tarde en 49 disponibles, la ventana es de 7: una tarde cada 7 semanas como
poco (figura~\ref{fig:rotacion}). Con la cuota del ejemplo anterior, 17 tardes en 48 semanas, la
ventana es de 2: nunca dos tardes seguidas. Como $g \ge 2$ siempre, nunca hay dos semanas
seguidas de la misma franja dura, pero sí puede haber una de tarde junto a una de noche.

\begin{figure}[ht]
\centering
\begin{tikzpicture}[x=0.62cm, y=0.62cm]
  \foreach \i in {1,...,21} {
    \pgfmathtruncatemacro{\t}{(\i==2 || \i==9 || \i==16) ? 1 : 0}
    \pgfmathtruncatemacro{\p}{(\i==5 || \i==12 || \i==19) ? 1 : 0}
    \ifnum\t=1
      \node[draw=black!50, fill=orange!45, minimum size=0.55cm, font=\scriptsize] at (\i,0) {T};
    \else\ifnum\p=1
      \node[draw=black!50, fill=green!30, minimum size=0.55cm, font=\scriptsize] at (\i,0) {P};
    \else
      \node[draw=black!50, fill=blue!12, minimum size=0.55cm, font=\scriptsize] at (\i,0) {M};
    \fi\fi
    \node[font=\tiny, black!60] at (\i,-0.75) {\i};
  }
  \draw[thick, decorate, decoration={brace, amplitude=5pt}] (5.6,0.55) -- (12.4,0.55)
    node[midway, above=6pt, font=\footnotesize] {cualquier ventana de 7 semanas: como mucho una T};
  \draw[thick, decorate, decoration={brace, amplitude=5pt, mirror}] (2,-1.2) -- (9,-1.2)
    node[midway, below=6pt, font=\footnotesize] {7 semanas entre dos tardes};
  \node[font=\footnotesize, anchor=west] at (21.8,0) {\dots};
\end{tikzpicture}
\caption{Rotación de un fijo con 7 semanas de tarde (T) en 49 disponibles: ventana de 7. Las
semanas de partido (P) tienen su propia ventana. Las mañanas (M) no tienen cuota ni ventana: son
lo que rellena entre medias.}
\label{fig:rotacion}
\end{figure}

En las franjas sin correturnos, donde la cuota aún no se conoce, se estima por arriba (las plazas
del año a cuatro días por semana, repartidas según $a_w$), para que la ventana no deje sin cubrir
lo que nadie más puede hacer.

\subsection{Objetivo lexicográfico}

Los criterios de la sección 1 no se mezclan en una suma ponderada: elegir pesos obliga a decidir
cuántas semanas de equidad vale una plaza sin cubrir, y una mala elección compra equidad con
cobertura. Se optimizan en orden estricto, por \emph{niveles}. Cada nivel elige, entre las
soluciones que dejó el anterior, las mejores según su criterio; el conjunto de soluciones se va
estrechando como en un embudo (figura~\ref{fig:niveles}).

\begin{figure}[ht]
\centering
\begin{tikzpicture}[
    niv/.style={draw, rounded corners=2pt, align=center, font=\small, minimum height=8mm},
    flecha/.style={-{Stealth[length=2mm]}, thick}
  ]
  \node[niv, minimum width=11cm, fill=gray!10] (r) {Cuadrantes que cumplen todas las restricciones};
  \node[niv, minimum width=9cm, fill=blue!8, below=6mm of r] (n1) {Niveles 0 y 1: los que cubren más plazas};
  \node[niv, minimum width=7cm, fill=blue!15, below=6mm of n1] (n2) {Niveles 2 y 3: los que dejan menos déficit de jornada};
  \node[niv, minimum width=5cm, fill=blue!25, below=6mm of n2] (n3) {Niveles 4 a 6: cuotas de semanas duras};
  \node[niv, minimum width=2.5cm, fill=blue!40, below=6mm of n3] (n4) {Cuadrante};
  \foreach \a/\b in {r/n1, n1/n2, n2/n3, n3/n4} \draw[flecha] (\a) -- (\b);
  \node[font=\footnotesize, align=left, anchor=west] at ([xshift=6mm]n2.east) {en cada paso: se clava\\el valor alcanzado y se\\siembra la solución};
\end{tikzpicture}
\caption{El objetivo por niveles. Cada nivel solo puede elegir entre las soluciones que dejan
intactos los valores de los niveles anteriores.}
\label{fig:niveles}
\end{figure}

\begin{center}
\begin{tabular}{@{}cll@{}}
  \toprule
  Nivel & Objetivo & Expresión \\
  \midrule
  0 & cobertura de lo que solo pueden hacer fijos & $\max \sum y_{w,d,k}$, casillas sin correturnos \\
  1 & cobertura & $\max \sum y_{w,d,k}$ \\
  2 & jornada: suma de déficits & $\min \sum_w \delta_w$ \\
  3 & jornada: peor déficit & $\min \max_w \delta_w$ \\
  4--6 & cuotas de noche, partido y tarde & $\min \sum_w \mathrm{dev}_{w,\varphi}$ \\
  \bottomrule
\end{tabular}
\end{center}
El nivel 0 va primero porque lo que no cubran los fijos lo cubren después los correturnos, salvo
las plazas a las que ellos no llegan (noches y guardias de otra base, líneas con cubridores). Si
los fijos no llegan a todo, que lo que quede sea lo que un correturno sí puede hacer. Los niveles
2 y 3 van juntos por la misma razón que en un reparto cualquiera: el primero acerca a todos a su
jornada; el segundo, entre las soluciones igual de buenas en total, evita que el déficit se
concentre en una sola persona.

Cada nivel $i$ se resuelve con su objetivo $f_i$; al terminar, su valor $v_i$ se \emph{clava} como
restricción ($f_i \ge v_i$ o $f_i \le v_i$) y se pasa al siguiente. Así ningún nivel puede
empeorar uno anterior. Cada nivel tiene un tiempo límite y, si no llega a demostrar el óptimo, se
clava el mejor valor encontrado: sigue siendo cierto que los niveles siguientes no lo empeoran.

Cada nivel arranca \emph{sembrado} con la solución del anterior: se le da como pista al solver.
No es solo una mejora de velocidad. Con los niveles anteriores clavados, encontrar \emph{alguna}
solución factible desde cero es muy difícil, y la del nivel anterior ya lo es por construcción.
Se siembran las variables de decisión y las auxiliares de los niveles ya clavados, pero no las
del nivel que empieza, que en la solución anterior valían cualquier cosa.

\subsection{Segunda fase: la línea concreta}

El modelo dice, por ejemplo, que $w$ trabaja de lunes a jueves en la casilla
(Valladolid, mañana, 480). Falta elegir la línea. Se hace sin solver, semana a semana y casilla a
casilla:
\begin{enumerate}[nosep]
  \item se atiende primero a quien tiene menos líneas posibles en común entre sus días;
  \item a cada uno se le da una línea que le valga \emph{todos} sus días de la semana (con plaza
        libre y C4 exacto contra lo que ya tiene), prefiriendo la de la semana anterior;
  \item si ninguna le vale la semana entera, se elige día a día, prefiriendo las que ya lleva esa
        semana.
\end{enumerate}
Así, siempre que se puede, una persona hace la misma línea toda la semana, y la misma que la
semana anterior si sigue en esa casilla. Como el modelo ya garantiza que el número de personas
por casilla cabe en sus plazas, esta fase casi nunca falla. Si queda algún día sin línea
compatible por C4, se informa. En los datos de 2027 no ha quedado ninguno. Antes de llegar a
esto se probó a resolver esta fase con un pequeño CP-SAT por semana: tardaba unos 150 segundos
frente a uno, y el resultado era peor en estabilidad.

% =========================================================================== %
\section{Correturnos y descansos}
% =========================================================================== %

\subsection{El modelo de los correturnos}

Lo que los fijos no cubren de lunes a viernes lo cubren los correturnos. El modelo es el mismo
que el de los fijos, con las mismas variables $y$ y las mismas restricciones: un turno al día,
no sobrecubrir, C4, C5, C6, descanso semanal, descanso de fin de semana y techo de jornada. Lo
que cambia responde a lo que es un correturno: la parte flexible de la plantilla.
\begin{itemize}[nosep]
  \item \textbf{Un solo modelo para toda la zona.} Un correturno trabaja en cualquier base (con
        la salvedad de noches y guardias de otra base, sección~\ref{sec:quien}), así que el
        problema ya no se separa por bases.
  \item \textbf{Cada línea es su propia casilla:}
        $\mathrm{cas}(l) = \bigl(b(l), \varphi(l), m(l), l\bigr)$. Como quedan pocas plazas, el
        modelo puede elegir la línea directamente, sin segunda fase, y C4 es exacto.
  \item \textbf{Sin zona semanal, sin cuotas y sin rotación.} Los correturnos absorben la
        inestabilidad que los fijos no tienen. La coherencia de su semana es solo una
        preferencia, el último nivel del objetivo.
  \item \textbf{No tienen que llegar a su jornada.} La jornada sigue siendo un techo, pero no
        hay déficit que minimizar: hay más horas de correturno que plazas por cubrir.
\end{itemize}
En la zona estudiada entran 8 correturnos, con 1.086 plazas pendientes y unas 7.750 variables.

\paragraph{Medida de desigualdad.} Para repartir algo con equidad entre $N$ personas con
recuentos $h_1, \dots, h_N$ y suma $S = \sum_v h_v$, se usa
\[
  \mathrm{Des}(h) \;=\; \sum_{w} \bigl| N h_w - S \bigr| \;+\; N\bigl(\max_w h_w - \min_w h_w\bigr).
\]
El primer término es $N$ veces la desviación de cada uno respecto a la media, escalada para no
usar números racionales; el segundo es el rango. Cada uno por separado es ciego a algo: con la
desviación sola, da igual que dos personas estén muy por encima o cuatro un poco por encima; con
el rango solo, únicamente cuentan los dos extremos. Ambos se linealizan con variables auxiliares
y las restricciones de máximo y mínimo del solver.

\paragraph{Objetivo.} Por niveles, igual que en los fijos:
\begin{center}
\begin{tabular}{@{}cll@{}}
  \toprule
  Nivel & Objetivo & Expresión \\
  \midrule
  1 & cobertura & $\max \sum y_{w,d,l}$ \\
  2 & equidad de horas & $\min \mathrm{Des}(h)$, con $h_w = M(w)$ \\
  3 & equidad de días duros & $\min \mathrm{Des}(u)$, con $u_w$ los días de tarde, partido o noche \\
  4 & coherencia semanal & $\min \sum_{w,s} \bigl( \sum_\zeta o_{w,s,\zeta} - t_{w,s} \bigr)$ \\
  \bottomrule
\end{tabular}
\end{center}
En el nivel 4, $o_{w,s,\zeta} = \max_{d \in s,\; \zeta(l) = \zeta} y_{w,d,l}$ indica si $w$ usa
la zona $\zeta$ (base y franja) en la semana $s$, y $t_{w,s} = \max_\zeta o_{w,s,\zeta}$ si
trabaja esa semana. El término cuenta las zonas que usa por encima de la primera: vale 0 si toda
la semana es de la misma base y franja. La equidad va antes que la coherencia porque, si
alguien tiene que cargar con la inestabilidad, que sea por igual.

\subsection{Descansos semanales}

Los dos modelos garantizan que cada semana completa tenga al menos dos días que pueden ser
descanso semanal, pero no dicen cuáles. El último paso los coloca, sin optimizar, para cada fijo
y cada correturno. Un día es candidato si está libre en el sentido de la sección 3: sin nada
asignado, sin vacaciones y sin ser un festivo entre semana.

Para cada semana completa $s$, hacen falta $\max(0,\, 2 - r_s)$ descansos, donde $r_s$ son los DS
y DO que ya trae el cuadrante (de un ciclo de patrón heredado). Se eligen así:
\begin{itemize}[nosep]
  \item si hacen falta dos, el \textbf{par de días libres seguidos más cercano al fin de semana}:
        dos días de descanso juntos valen más que dos sueltos, y pegados al fin de semana más
        todavía;
  \item si no hay ningún par, o solo hace falta uno, los últimos días libres de la semana.
\end{itemize}
En las semanas partidas por vacaciones o por el principio o el final del año no se exige el
descanso semanal: se ponen hasta dos con los días libres que queden.

Como los modelos ya reservan el sitio, este paso solo puede fallar en una semana completa si lo
que traían los pasos anteriores ya no dejaba dos días libres; en ese caso los modelos lo toleran
y no añaden más trabajo esa semana. En la ejecución con los datos de 2027 no ha quedado ningún
descanso sin colocar.

% =========================================================================== %
\section{Resultados y limitaciones}
% =========================================================================== %

\subsection{El caso estudiado}

Los resultados corresponden a la zona de Valladolid para el año 2027:
\begin{itemize}[nosep]
  \item \textbf{6 bases}: Valladolid, Medina, Íscar, Mayorga, Tordesillas y Peñafiel;
  \item \textbf{71 líneas}: 40 de mañana, 19 de tarde, 6 de partido, 3 de noche y 3 guardias
        localizadas;
  \item \textbf{87 trabajadores}: 8 de patrón (4 patrones: dos de noches y dos de UVI), 71 fijos y
        8 correturnos;
  \item jornada de 1776~h, descanso mínimo de 12~h, 6 días y 48~h por semana como máximo.
\end{itemize}
Se ha ejecutado con un tiempo límite de 600~s por nivel de cada modelo y 8 hilos. La ejecución
completa tarda unos 70 minutos; casi todo ese tiempo se va en los niveles que agotan su límite
sin demostrar el óptimo.

\subsection{Cobertura}

\begin{center}
\begin{tabular}{@{}lrrr@{}}
  \toprule
  Paso & Plazas & Cubiertas & Sin cubrir \\
  \midrule
  Festivos, sábados y domingos (voraz) & --- & todas & 0 \\
  Lunes a viernes: fijos & 15.189 & 14.122 & 1.067 \\
  Lunes a viernes: correturnos & 1.067 & 995 & 72 \\
  \midrule
  Total de lunes a viernes & 15.189 & 15.117 & \textbf{72} (0,5\,\%) \\
  \bottomrule
\end{tabular}
\end{center}
Las plazas de lunes a viernes son las que quedan tras los patrones y sus libranzas. Las 72 que
quedan sin cubrir son sobre todo de mañana:
\begin{center}
\begin{tabular}{@{}lrrr@{}}
  \toprule
  Base & Mañana & Tarde & Partido \\
  \midrule
  Valladolid & 18 & 15 & 9 \\
  Íscar & 9 & 3 & --- \\
  Tordesillas & 3 & 7 & --- \\
  Medina & 5 & --- & --- \\
  Mayorga & 3 & --- & --- \\
  \bottomrule
\end{tabular}
\end{center}
La cobertura del modelo de correturnos es óptima demostrada: con lo que dejan los pasos
anteriores, no se puede cubrir ni una plaza más sin incumplir alguna restricción. Los fijos ya
están en su jornada, así que esas plazas solo se cubrirían con más plantilla o cambiando
decisiones de pasos anteriores.

\subsection{Legalidad}

El cuadrante final no tiene ningún incumplimiento de C4, C5 ni C6 en lo que decide el
algoritmo, y todos los descansos semanales caben: ninguna semana completa se queda sin sus dos
DS. Las únicas formas fuera del convenio son las de los patrones pactados, que se dan por
válidas (sección 3).

\subsection{Jornada}

\begin{itemize}[nosep]
  \item \textbf{Fijos}: 70 de 71 terminan exactamente en su jornada (mediana y máximo de
        1776~h). El único que se queda por debajo, con 1744~h, es un cubridor de UVI, cuyo año
        está condicionado por los ciclos que hereda.
  \item \textbf{Correturnos}: entre 1200 y 1208~h. La diferencia entre ellos es de un turno. No
        llegan a su jornada porque no hay trabajo para ello: quedan unas 570~h de cada uno sin
        usar, que es la holgura de la plantilla.
\end{itemize}

\subsection{Equidad}

\paragraph{Fines de semana y festivos.} Recuento por persona dentro de cada base:
\begin{center}
\begin{tabular}{@{}lrccc@{}}
  \toprule
  Base & Personas & Sábados & Domingos & Festivos \\
  \midrule
  Valladolid (con correturnos) & 53 & 15--16 & 6--7 & 3--5 \\
  Medina & 8 & 12--13 & 6--7 & 3--4 \\
  Íscar & 3 & 16--17 & --- & 4--5 \\
  Tordesillas & 3 & 16--17 & --- & 4--5 \\
  Mayorga & 4 & 12--13 & --- & 3--4 \\
  \bottomrule
\end{tabular}
\end{center}
La diferencia es de un día como mucho, salvo en los festivos de Valladolid (dos). Además quedan
espaciados: nadie trabaja más de 3 fines de semana seguidos, y el hueco más largo entre dos
fines de semana trabajados es de 9 semanas (6,3 de media).

\paragraph{Semanas duras.} De los 58 fijos con cuota de semanas de tarde, 54 la cumplen
exactamente, 3 se quedan a una semana y 1 a dos. En las noches de Medina, 5 de 8 la cumplen
exactamente. Entre los correturnos, los días duros son entre el 43\,\% y el 45\,\% de sus días
de lunes a viernes.

\subsection{Rotación y estabilidad}

\begin{itemize}[nosep]
  \item Ningún fijo encadena dos semanas seguidas de la misma franja dura, y ninguna separación
        entre dos semanas de una misma franja queda por debajo de su ventana (922 separaciones).
  \item Sí hay 43 pares de semanas duras de franjas distintas seguidas (una de tarde junto a
        una de noche, por ejemplo). Es lo que permite la restricción, y lo que se decidió.
  \item Los correturnos trabajan en una sola zona (base y franja) entre 38 y 42 de sus unas 48
        semanas con trabajo.
\end{itemize}

\subsection{Comportamiento del solver}

Los modelos de las bases pequeñas (Íscar, Mayorga, Tordesillas y Peñafiel) se resuelven al
óptimo demostrado en todos sus niveles en segundos. En Medina, la cobertura y la jornada llegan
al óptimo; las cuotas, no. En Valladolid, la cobertura agota los 600~s sin demostrar el óptimo
y se queda en el mismo valor que con 60~s, señal de que está muy cerca de su máximo; la jornada
sí se demuestra óptima. En los correturnos, la cobertura es óptima en un segundo, y la
equidad de horas y la coherencia semanal se quedan en factible.

Que un nivel no demuestre su óptimo no rompe el orden de prioridades: el valor encontrado se
clava igual y ningún nivel posterior lo empeora. Lo que se pierde es la garantía de que no
existía algo mejor en ese nivel.

\subsection{Limitaciones y trabajo futuro}

\begin{itemize}
  \item \textbf{Una sola zona.} El diseño separa convenio, zona y datos, pero solo existe la
        receta de Valladolid. Cada zona nueva necesita la suya, y es probable que obligue a
        sacar a la parte común piezas que hoy son de Valladolid.
  \item \textbf{Datos provisionales.} La carga de datos lee ficheros CSV preparados a mano. Su
        formato cambiará cuando se conecte con las fuentes reales de la empresa, y el resto del
        algoritmo no debería verse afectado.
  \item \textbf{Restricciones no modeladas.} Las conciliaciones admiten tipos de día, líneas y
        ventanas horarias, pero no reglas del tipo \emph{``como mucho un domingo al mes''}.
  \item \textbf{Huecos agrupados.} Las plazas que quedan sin cubrir están dispersas. Una mejora
        con valor práctico es un nivel que, sin perder cobertura, las concentre en las mismas
        líneas y semanas seguidas, de modo que se puedan cubrir con contratos temporales en
        lugar de días sueltos.
  \item \textbf{Optimalidad.} Los niveles grandes no demuestran su óptimo en el tiempo dado. El
        problema se ejecuta una vez al año, así que se prefiere dar más tiempo a recortar
        niveles; aun así, la descomposición en pasos hace que el resultado sea óptimo dentro de
        cada paso, no para el año completo.
\end{itemize}

\end{document}
